\documentclass[11pt,letterpaper]{article}
\usepackage[margin=0.9in]{geometry}
\usepackage{amsmath,amssymb,bm,booktabs,longtable,array,xcolor,hyperref,microtype}
\hypersetup{colorlinks=true,linkcolor=blue!45!black,urlcolor=blue!45!black,
 pdftitle={Analytical gradients of determinantal CC(P;Q) and EOMCC(P;Q)},
 pdfauthor={CCpy development notes}}
\setlength{\parindent}{0pt}
\setlength{\parskip}{5pt}
\setlength{\emergencystretch}{2em}
\allowdisplaybreaks[2]
\numberwithin{equation}{section}
\renewcommand{\arraystretch}{1.13}
\newcommand{\ket}[1]{\lvert #1\rangle}
\newcommand{\bra}[1]{\langle #1\rvert}
\newcommand{\me}[3]{\langle #1\rvert #2\lvert #3\rangle}
\newcommand{\bH}{\overline H}
\newcommand{\T}{\mathsf T}
\newcommand{\dd}{\mathrm d}
\newcommand{\ee}[2]{#1\times10^{#2}}
\newcommand{\eunit}{E_{\mathrm h}}
\newcommand{\gunit}{E_{\mathrm h}/a_0}
\newcommand{\vecz}{\mathbf z}
\newcommand{\evec}{\mathbf e_0}
\newcommand{\Pzero}{\mathbb P_0}
\newcommand{\Qproj}{\mathbb Q}
\newcommand{\sym}{\operatorname{sym}}
\newcommand{\diag}{\operatorname{diag}}
\newcommand{\supp}{\operatorname{supp}}
\newcommand{\tr}{\operatorname{Tr}}
\title{Analytical gradients of determinantal\\CC(P;Q) and EOMCC(P;Q)\\[4pt]
\large Stationary functionals, response equations, and validation in CCpy}
\author{Implementation-linked theory note}
\date{2 November 2026}
\begin{document}
\maketitle
\begin{abstract}
We derive first nuclear derivatives of a specified biorthogonal, noniterative
CC(P;Q) energy for fixed, general determinant selections. Ground and excited
states share the moment correction, but their defining constraints and
response variables differ. Two separate stationary functionals produce the
base-state and correction derivatives. The derivation retains the distinction
between the CC residual Jacobian and the projected similarity-transformed
Hamiltonian, the distinction between a projected left state and a ground-state
gradient multiplier, the excited right-reference coefficient, and the explicit
bare-Fock denominator derivative. We pull the stationary coefficients back to
bare-integral density matrices, derive canonical-RHF orbital relaxation and
AO-overlap terms, and give matching density-fitting/Cholesky contractions.
Every construction is mapped to the determinant-driven implementation. Tables
are generated from saved numerical reports, including independent optimized
CCpy energy finite differences. The resulting evidence validates derivatives
of the stated models and tested spaces; it is neither a proof of physical
energy accuracy nor a validation of every approximation bearing a CC(P;Q) name.
\end{abstract}
\bigskip
\begin{center}\textbf{The stationary construction}\end{center}
\[
 \boxed{\mathcal L_a=F_a+\mathbf z_a^{\T}\mathbf c,
 \qquad c_y^{\T}\mathbf z_a=-F_{a,y},\qquad a=b,c.}
\]
The base energy and Q correction have separate amplitude adjoints. Their
bare-integral derivative densities are then relaxed with respect to the
reference orbitals before contraction with nuclear derivative integrals.
\clearpage
{\small\setlength{\baselineskip}{10pt}\setlength{\parskip}{0pt}\tableofcontents}
\clearpage

\section{Scope, conventions, and the energy being differentiated}
\label{sec:definitions}
We use real orbitals, real isolated roots, and a finite orthonormal
spin-orbital basis. Transposes below are algebraic transposes of the real
biorthogonal problem, not assertions that the similarity-transformed
Hamiltonian is Hermitian. Nuclear repulsion is excluded from the electronic
Hamiltonian until the final molecular gradient. The implemented molecular
reference is closed-shell canonical RHF; excited states can be singlets or
triplets in its fixed-$M_s=0$ determinant sector. An open-shell reference
would require its own orbital-response derivation. The molecular derivation
assumes a linearly independent AO basis with its full corresponding MO span
retained, and no changing linear-dependence cutoff.

Let $\ket{0}$ be the reference determinant. Let $P$ and $Q$ be fixed, disjoint
sets of excited determinants, neither containing $0$, and define
\begin{equation}
 \Pzero=\ket{0}\bra{0}+\sum_{I\in P}\ket{I}\bra{I},\qquad
 \Qproj=\sum_{K\in Q}\ket{K}\bra{K}.
 \label{eq:projectors}
\end{equation}
$Q$ need not be the entire complement of $P$ in the full Hilbert space.
For example, it can be the remaining triples, or remaining triples and
quadruples. Omitted determinants outside both $P$ and $Q$ are not silently
included in the correction. The full determinant support used for actions
of $H$ and $e^{\pm T}$ is not restricted to $\{0\}\cup P\cup Q$.

We choose pure particle-hole excitation operators with phases fixed by
$\tau_I\ket{0}=\ket{I}$. In this unique-determinant convention,
\begin{equation}
 T=\sum_{I\in P}t_I\tau_I,\qquad U=e^{-T},\quad V=e^T,\quad
 \bH=UHV.
 \label{eq:similarity}
\end{equation}
There are no factorial multiplicities in this sum over unique excitations.
Conversion to redundant antisymmetric CCpy spin blocks must restore their
operator phases and permutation images. Because pure reference excitations
commute with one another,
\begin{equation}
 \frac{\partial\bH}{\partial t_I}
   =\bH\tau_I-\tau_I\bH\equiv H^{I}.
 \label{eq:commutator}
\end{equation}
The exponentials are finite on the determinant support by nilpotency. This
identity does not require $P$ to be closed under de-excitation.

\subsection{The base states and correction}
Write the uncorrected electronic state energy as $E_\mu$, with $\mu=0$ for
the CC(P) ground branch. The corrected energy is
\begin{equation}
 \mathcal E_\mu=V_{\rm nn}+E_\mu+\delta_\mu,\qquad
 \delta_\mu=\sum_{K\in Q}\frac{N_{\mu K}M_{\mu K}}{D_{\mu K}},
 \label{eq:energy}
\end{equation}
where
\begin{equation}
 M_{\mu K}=\me{K}{\bH}{R_\mu},\qquad
 N_{\mu K}=\me{L_\mu}{\bH}{K},\qquad
 \langle L_\mu\vert R_\mu\rangle=1.
 \label{eq:moments}
\end{equation}
Here $\ket{R_\mu}$ and $\bra{L_\mu}$ are vectors in the
\emph{similarity-transformed} reference+$P$ space. The corresponding full
CC wavefunctions are $e^T\ket{R_\mu}$ and $\bra{L_\mu}e^{-T}$.
Without normalization, the right-hand side of the correction would have to
be divided by $\langle L_\mu\vert R_\mu\rangle$.

One way to motivate this energy, including the choice of left state, is to
hold T(P) fixed and introduce a coupling parameter $\eta$ into the retained reference+P/Q block:
\begin{equation}
 \begin{pmatrix}\mathsf A&\eta\bH_{\{0,P\},Q}\\
 \eta\bH_{Q,\{0,P\}}&\bH_{QQ}\end{pmatrix}
 \begin{pmatrix}\bm R_P\\\bm R_Q\end{pmatrix}
 =\mathcal E^{\rm el}(\eta)\begin{pmatrix}\bm R_P\\\bm R_Q\end{pmatrix}.
\end{equation}
The Q equation gives
\begin{equation}
 \bm R_Q=\eta[\mathcal E^{\rm el}(\eta)\mathbf1-\bH_{QQ}]^{-1}
                         \bH_{Q,\{0,P\}}\bm R_P.
\end{equation}
Substitute into the P equation and project with the unperturbed left
eigenvector. For a simple unperturbed eigenvalue and $\bm l^\T\bm r=1$,
\begin{equation}
 \mathcal E^{\rm el}(\eta)=E_\mu+\eta^2
 \bm l^\T\bH_{\{0,P\},Q}(E_\mu\mathbf1-\bH_{QQ})^{-1}
 \bH_{Q,\{0,P\}}\bm r+O(\eta^4).
 \label{eq:resolventorigin}
\end{equation}
A diagonal Q-resolvent model gives Eq.~\eqref{eq:energy}. If Q omits other
determinants, this is an elimination in the retained block, not an exact
elimination of the complete Hilbert-space complement. The artificial
parameter organizes this block coupling; it is not a claim of a small
physical expansion parameter at $\eta=1$. This construction selects the
\emph{projected left eigenstate} for N, not the ground CC residual-gradient
multiplier introduced below. Replacing N by a conjugated right moment
would also be unjustified for this non-Hermitian problem.

The two implemented denominator models are
\begin{align}
 D_{\mu K}^{\rm full}&=E_\mu-\bH_{KK}, \label{eq:denfull}\\
 D_{\mu K}^{\rm mp}&=\omega_\mu+\sum_p\chi_{Kp}f_{pp},\qquad
 \omega_\mu=E_\mu-E_0,\qquad
 \chi_{Kp}=n_p^0-n_p^K. \label{eq:denmp}
\end{align}
The Fock operator is the \emph{bare reference Fock operator}, not a
similarity-transformed one. Its occupied sum includes frozen core electrons.
The MP model can be written as $D_{\mu K}=E_\mu-E_K^{\rm mp}$ with
$E_K^{\rm mp}=E_0-\sum_p\chi_{Kp}f_{pp}$, namely the CC(P) ground energy
plus the reference orbital-energy excitation gap for K.
For the ground state, $\omega_0=0$. In the optimized SD/active-SDt comparisons
reported here, \texttt{mp} matches variant A and \texttt{full} matches variant D.
This naming correspondence is not a definition of every published A--D
approximation for arbitrary excitation ranks.

Equation~\eqref{eq:energy} is noniterative: $E_\mu$, $T$, $R_\mu$ and $L_\mu$
are obtained from the \emph{uncorrected} P-space problem. One does not replace
$E_\mu$ in a denominator or eigenvalue equation by $E_\mu+\delta_\mu$.
Such a replacement defines a different implicit energy and a different
response problem. In particular, no extra universal normalization factor
$1/(1+\sum_K NM/D^2)$ belongs to the model in Eq.~\eqref{eq:energy}.

The supplied ground-state derivation is the internal note
\path{ccpy/gradients/ccpq_gradients_revised.pdf} (implementation revision
1 November 2026). The \texttt{ccpq\_audit} extension supplies the unified
projected-state audit. The present note gives the ground/excited derivation
and the subsequent molecular/optimized-code evidence explicitly; it does not
rely on the upstream audit as a claim of verified nuclear gradients.

\subsection{Bare versus normal-ordered energies}
If an implementation uses $H=E_{\rm ref}^{\rm el}+H_N$, then
$\bH=E_{\rm ref}^{\rm el}\mathbf 1+\overline H_N$. Off-diagonal P--Q moments
and both denominator models are unchanged by this scalar shift, while
\begin{equation}
 E_\mu=E_{\rm ref}^{\rm el}+E_{\mu,N},\qquad
 \mathcal E_\mu=V_{\rm nn}+E_{\rm ref}^{\rm el}+E_{\mu,N}+\delta_\mu.
\end{equation}
The determinant code uses the bare electronic Hamiltonian, so its densities
already contain the reference and frozen-core contributions. Adding an HF
reference derivative to these densities a second time would double count it.
The optimized tensor comparison reconstructs the same total state energies
from reference energy, CC correlation energy, excitation energy and correction.

\section{Right CC(P), the two ground left quantities, and the Jacobian}
The amplitudes solve
\begin{equation}
 m_A(t,H)=\bH_{A0}=0\quad(A\in P),\qquad E_0=\bH_{00}.
 \label{eq:ccright}
\end{equation}
Let
\begin{equation}
 e_I=\frac{\partial E_0}{\partial t_I}=\bH_{0I},\qquad
 J_{AI}=\frac{\partial m_A}{\partial t_I}
       =\me{A}{[\bH,\tau_I]}{0},\qquad
 B_{AI}=\bH_{AI}-E_0\delta_{AI}.
 \label{eq:JBe}
\end{equation}
To see exactly why $J$ and $B$ can differ, expand
\begin{align}
 J_{AI}&=\bH_{AI}-\sum_X\me{A}{\tau_I}{X}\bH_{X0}\\
 &=B_{AI}-\sum_{X\notin\{0\}\cup P}
                \me{A}{\tau_I}{X}\bH_{X0}
       \equiv B_{AI}-\mathcal W_{AI}.
 \label{eq:JminusB}
\end{align}
The $P$ residuals vanish, but omitted residuals generally do not. A selected
high-rank determinant can receive an excitation of an omitted lower-rank
determinant. The last term then survives. De-excitation closure of $P$ is a
sufficient condition for $\mathcal W=0$. In particular, all S/D plus selected
triples has this property: an omitted residual has rank at least three and
cannot be raised into a P determinant of rank at most three.

The projected left ground state used in the numerator is
\begin{equation}
 \bra{L_0}=\bra{0}+\sum_{A\in P}\lambda_A\bra{A},\qquad
 B^{\T}\bm\lambda=-\bm e,\qquad \ket{R_0}=\ket{0}.
 \label{eq:projectedlambda}
\end{equation}
The base-energy \emph{gradient multiplier} is instead
\begin{equation}
 J^{\T}\bm\lambda^{g}=-\bm e,\qquad
 \mathcal L_b^{(0)}=E_0+(\bm\lambda^{g})^{\T}\bm m.
 \label{eq:lambdagrad}
\end{equation}
It follows that $\partial_t\mathcal L_b^{(0)}=0$ and
$\partial_{\lambda^g}\mathcal L_b^{(0)}=\bm m=0$.
Using the projected $\bm\lambda$ in its place gives the defect
\begin{equation}
 \bm e+J^{\T}\bm\lambda=(J-B)^{\T}\bm\lambda
                         =-\mathcal W^{\T}\bm\lambda.
 \label{eq:lambdadefect}
\end{equation}
Thus $\lambda=\lambda^g$ is a consequence of additional structure, not a
valid general assumption. The program retains both quantities. Its
nonclosed-P tests explicitly verify nonzero differences and the stationary
base derivative with $\lambda^g$.

\section{Projected excited states and their reference component}
Let $\mathsf A$ be the reference+$P$ projection of $\bH$. At the CC solution,
\begin{equation}
 \mathsf A=\begin{pmatrix}E_0&\bm h^{\T}\\0&\bH_{PP}\end{pmatrix},
 \quad h_I=\bH_{0I},\qquad
 \mathsf A\bm r_\mu=E_\mu\bm r_\mu,\quad
 \mathsf A^{\T}\bm l_\mu=E_\mu\bm l_\mu.
 \label{eq:eom}
\end{equation}
For $E_\mu\ne E_0$,
\begin{equation}
 l_{\mu0}=0,\qquad
 r_{\mu0}=\frac{\bm h^{\T}\bm r_{\mu P}}{E_\mu-E_0},\qquad
 \bm l_{\mu P}^{\T}\bm r_{\mu P}=1.
 \label{eq:r0}
\end{equation}
The left reference coefficient vanishes; the right one generally does not.
Consequently,
\begin{equation}
 M_{\mu K}=r_{\mu0}\bH_{K0}+\sum_{I\in P}\bH_{KI}r_{\mu I},\qquad
 N_{\mu K}=\sum_{I\in P}l_{\mu I}\bH_{IK}.
 \label{eq:r0moment}
\end{equation}
Dropping $r_{\mu0}$ changes the correction whenever its first term is
nonzero. The H4 singlet benchmark has $|r_{\mu0}|\simeq0.083$, so that check
is not restricted to a symmetry-forced zero-reference case.

For a general selected P, Eq.~\eqref{eq:eom} is the defining projected-Hbar
EOM problem in this note. The CC residual Jacobian $J$ is still used for
\emph{cluster response}. Substituting $J$ for the projected EOM matrix in
an arbitrary nonclosed P changes the chosen base-state model and moments.
For all-S/D plus selected-T spaces, the usual corresponding optimized EOM
problem agrees with this projected definition, as verified independently.

\section{A stationary representation of the Q correction}
Suppress $\mu$ temporarily and define divided moments
\begin{equation}
 u_K=\frac{M_K}{D_K},\qquad v_K=\frac{N_K}{D_K},\qquad
 w_K=u_Kv_K=\frac{N_KM_K}{D_K^2},\qquad W_Q=\sum_{K\in Q}w_K.
 \label{eq:divided}
\end{equation}
The auxiliary saddle functional
\begin{equation}
 \mathcal C_Q=\sum_{K\in Q}\bigl(v_KM_K+u_KN_K-u_Kv_KD_K\bigr)
 \label{eq:saddleQ}
\end{equation}
has derivatives
\begin{equation}
 \partial_{u_K}\mathcal C_Q=N_K-v_KD_K=0,\qquad
 \partial_{v_K}\mathcal C_Q=M_K-u_KD_K=0.
\end{equation}
At these solutions, $\mathcal C_Q=\delta$. The envelope derivative is
\begin{equation}
 \dd\delta=\sum_K\bigl(v_K\dd M_K+u_K\dd N_K-w_K\dd D_K\bigr).
 \label{eq:deltadifferential}
\end{equation}
This is also obtained directly by differentiating $NM/D$; the saddle form
explains why no $u^x$ or $v^x$ is required in a first derivative.
On shell its three terms give $\delta+\delta-\delta$. Omitting the last
term gives twice the correction and loses the denominator derivative.
It is the Q-sector subtraction associated with $(\bH-E_\mu)$, not an
optional bookkeeping constant.

The denominator differentials are
\begin{align}
 \dd D_{\mu K}^{\rm full}&=\dd E_\mu-\dd\bH_{KK},\\
 \dd D_{\mu K}^{\rm mp}&=\dd E_\mu-\dd E_0+
                                     \sum_p\chi_{Kp}\dd f_{pp}.
 \label{eq:dendifferentials}
\end{align}
In particular the full-diagonal term contributes $+w_K\dd\bH_{KK}$,
whereas the bare-Fock contribution is $-w_K\sum_p\chi_{Kp}\dd f_{pp}$.
The sign follows from differentiating $1/D$; it is not a convention to be
chosen independently.

For example the explicit total differential before eliminating responses is
\begin{align}
 \dd(E_\mu+\delta_\mu)_{\rm full}
   &=(1-W_Q)\dd E_\mu+
          \sum_K\bigl(v_K\dd M_K+u_K\dd N_K+w_K\dd\bH_{KK}\bigr),\\
 \dd(E_\mu+\delta_\mu)_{\rm mp}
   &=(1-W_Q)\dd E_\mu+W_Q\dd E_0+
      \sum_K(v_K\dd M_K+u_K\dd N_K)
      -\sum_{Kp}w_K\chi_{Kp}\dd f_{pp}.
 \label{eq:chainbeforeadjoint}
\end{align}
For $\mu=0$ the first two MP energy terms combine to $\dd E_0$.
These equations still contain implicit T/R/L response. They are not yet
AO nuclear-gradient formulas, and do not justify omitting those responses.

\section{The general constrained-adjoint construction}
\label{sec:adjoint}
Let $\bm y$ contain the independently varied defining variables and let
$\bm c(\bm y,H)=0$ be their constraints. For either $a=b$ (base) or
$a=c$ (correction), define
\begin{equation}
 \boxed{\mathcal L_a=F_a(\bm y,H,f)+\vecz_a^{\T}\bm c(\bm y,H),\qquad
 c_y^{\T}\vecz_a=-F_{a,y}.}
 \label{eq:masterfunctional}
\end{equation}
The Fock diagonal is displayed separately because MP has an explicit Fock
dependence. It is ultimately a function of the same bare integrals and
reference occupation, not an unrelated external parameter.

To derive the gradient, differentiate the constraints and objective:
\begin{align}
 c_y\bm y^x+c_x&=0,\qquad \bm y^x=-c_y^{-1}c_x,\\
 \frac{\dd F_a}{\dd x}
 &=F_{a,x}+F_{a,y}^{\T}\bm y^x
   =F_{a,x}-F_{a,y}^{\T}c_y^{-1}c_x
   =F_{a,x}+\vecz_a^{\T}c_x
   =\left.\partial_x\mathcal L_a\right|_{\bm y,\vecz_a}.
 \label{eq:adjointproof}
\end{align}
Thus all defining variables are held fixed in the final explicit integral
derivative. The total molecular functional is
\begin{equation}
 \mathcal L_{\rm tot}=V_{\rm nn}+\mathcal L_b+\mathcal L_c.
 \label{eq:separatefunctionals}
\end{equation}
The two adjoints can be added because their constraint Jacobian is the same.
Keeping them separate yields independent base/correction densities, makes
sign errors visible, and is exactly how the determinant driver reports the
three gradients. A single stationary functional for $F_b+F_c$ is equivalent
if all sources and constraints are retained.
Separate stationarity defines the individual multipliers; stationarity of
the sum alone fixes only the combined adjoint, not its base/correction split.

An isolated root, a valid gauge, nonzero Q denominators and nonsingular
constraint response are assumed. A root crossing, a changing discrete P/Q
selection, or a genuinely singular reference Hessian requires additional
analysis; it is not repaired by writing an adjoint equation.

\section{Expanded ground-state response: Sigma and Pi}
\label{sec:ground}
Take $\bm y=(\bm t,\bm\lambda)$, $F_b=E_0$, $F_c=\delta_0$ and
\begin{equation}
 \bm c_0=\begin{pmatrix}\bm m\\B^{\T}\bm\lambda+\bm e\end{pmatrix}.
 \label{eq:groundconstraints}
\end{equation}
The base adjoint has zero multiplier on the left-state constraint and
$\lambda^g$ on the right CC equations, recovering Eq.~\eqref{eq:lambdagrad}.
Let the correction multipliers be $\bm\pi$ and $\bm\sigma$. Then
\begin{equation}
 \boxed{\mathcal L_c^{(0)}=\mathcal C_Q+\bm\pi^{\T}\bm m+
       \bra{L_0}(\bH-E_0)\ket{\Sigma},\qquad
       \ket{\Sigma}=\sum_{I\in P}\sigma_I\ket{I}.}
 \label{eq:groundcorrectfunctional}
\end{equation}
Indeed its last term is
$\bm\sigma^{\T}(B^{\T}\bm\lambda+\bm e)$, which vanishes on the projected
left equations. The functional therefore equals $\delta_0$ on shell.

\subsection{Left-state response}
Only $N_K$ and the constraint term depend explicitly on $\lambda_I$.
Using $\partial_{\lambda_I}N_K=\bH_{IK}$ gives
\begin{equation}
 (s_\lambda)_I=\sum_{K\in Q}\bH_{IK}u_K,
 \qquad \boxed{B\bm\sigma=-\bm s_\lambda.}
 \label{eq:sigma}
\end{equation}
It is $B$, not $J$, in this equation: the constraint being differentiated
is the projected-left-state equation.

\subsection{Cluster response}
Hold $\lambda,u,v,\sigma$ fixed and set
\begin{align}
 (s_t)_I&=\sum_K\left[v_KH^I_{K0}
                   +u_K\me{L_0}{H^I}{K}-w_K d_{KI}\right],\\
 d_{KI}^{\rm full}&=e_I-H^I_{KK},\qquad d_{KI}^{\rm mp}=0,\\
 k_I&=\me{L_0}{H^I}{\Sigma}
                      -(\bm\lambda^{\T}\bm\sigma)e_I.
 \label{eq:groundsource}
\end{align}
The MP Fock diagonal is independent of T at fixed bare integrals/reference;
its separate integral derivative enters later. Stationarity in $t_I$ yields
\begin{equation}
 \boxed{J^{\T}\bm\pi=-\bm s_t-\bm k.}
 \label{eq:pi}
\end{equation}
The sequence is consequently: right CC(P), projected left state,
$\lambda^g$, Q divided moments, $\sigma$, then $\pi$. Solving the Sigma
equation with $J$ or the Pi equation with $B^{\T}$ is generally incorrect.

\subsection{The ground Hbar coefficient}
Embed all vectors in the full determinant support. Write
$\bm b=\evec+\bm\lambda_P$, $\bm u_Q$, $\bm v_Q$, $\bm\sigma_P$ and
$\bm\pi_P$ for those embeddings. If
$\mathcal A_{AB}=\partial\mathcal L/\partial\bH_{AB}$ at fixed multipliers,
then
\begin{align}
 \mathcal A_b^{(0)}&=(\evec+\bm\lambda_P^g)\evec^{\T},\\
 \mathcal A_c^{(0)}&=\bm\pi_P\evec^{\T}
   +\bm b\bm\sigma_P^{\T}
   -(\bm\lambda^{\T}\bm\sigma)\evec\evec^{\T}
   +\bm v_Q\evec^{\T}+\bm b\bm u_Q^{\T}+\mathcal D_0,\\
 \mathcal D_0^{\rm full}&=\sum_{K\in Q}w_K\ket{K}\bra{K}
                                     -W_Q\ket{0}\bra{0},\qquad
 \mathcal D_0^{\rm mp}=0.
 \label{eq:groundcoeff}
\end{align}
For MP, this coefficient must be supplemented by the explicit Fock density
in Sec.~\ref{sec:fock}. For full denominators every denominator contribution
is already in $\mathcal D_0$.

\section{Expanded excited-state constrained response}
\label{sec:excited}
Let $k=1+|P|$. Fix a nonzero component $r_q=c$ of the right vector; $q$ is
chosen from P and $c$ is constant during differentiation. Let $D_q$ delete
component $q$ of a vector and let $E_q$ insert a zero at that component.
The independent variables and constraints are
\begin{align}
 \bm y&=(\bm t,\bm r_{\widehat q},\bm l,E),\\
 \bm c_\mu&=\begin{pmatrix}
  \bm m\\(\mathsf A-E\mathbf1)\bm r\\
  D_q(\mathsf A^{\T}-E\mathbf1)\bm l\\\bm l^{\T}\bm r-1
 \end{pmatrix}.
 \label{eq:excconstraints}
\end{align}
There are $|P|+2k$ variables and constraints. One left eigen-equation is
removed because, given the right equations, the two sets have one algebraic
redundancy. Fixing the right scale and imposing $\bm l^{\T}\bm r=1$
removes the eigenvector scaling freedom. Paired changes
$\bm r\mapsto c'\bm r$, $\bm l\mapsto\bm l/c'$ leave the energy and densities
unchanged, including negative real scale changes on a regular gauge patch.

The objectives used by the unified oracle are
\begin{equation}
 F_b=\bm l^{\T}\mathsf A\bm r,\qquad F_c=\delta_\mu.
 \label{eq:objectives}
\end{equation}
On shell $F_b=E$. Off shell it is important to use derivatives of the
chosen expression consistently. Replacing $F_b$ by $E$ is possible, but
changes the adjoints because
\begin{equation}
 \bm l^{\T}\mathsf A\bm r-E
 =\bm l^{\T}(\mathsf A-E\mathbf1)\bm r
                    +E(\bm l^{\T}\bm r-1).
 \label{eq:offshellequivalence}
\end{equation}
The difference is a combination of constraints, not an identity before
those constraints are imposed.

\subsection{State and energy sources}
At fixed $t$, use full reference+$P$ components for $r,l$. The base sources are
\begin{equation}
 F_{b,r}=\mathsf A^{\T}\bm l,\quad F_{b,l}=\mathsf A\bm r,
 \quad F_{b,E}=0,\quad (F_{b,t})_I=\me{L_\mu}{H^I}{R_\mu}.
 \label{eq:baseeomsources}
\end{equation}
For the correction,
\begin{align}
 F_{c,r}&=\bH_{Q,\{0,P\}}^{\T}\bm v,
 &F_{c,l}&=\bH_{\{0,P\},Q}\bm u,
 &F_{c,E}&=-W_Q,\\
 (F_{c,t})_I&=\sum_K\left[v_K\me{K}{H^I}{R_\mu}
       +u_K\me{L_\mu}{H^I}{K}\right]+d_I,\\
 d_I^{\rm full}&=\sum_Kw_K H^I_{KK},
 &d_I^{\rm mp}&=W_Q e_I.
 \label{eq:excsources}
\end{align}
The $E$ in these partial derivatives is independent. In particular,
$\partial_t D^{\rm mp}=-\partial_t E_0$ and
$\partial_E D^{\rm mp}=1$. Collapsing $E$ into a T-dependent expression
before differentiating these sources, while also retaining its response
constraint, would double count its chain rule.

\subsection{Explicit adjoint equations and block order}
Use multipliers $\bm z_t$, $\bm a\in\mathbb R^k$,
$\bm b\in\mathbb R^k$ with $b_q=0$, and $\nu$. The functional is
\begin{align}
 \mathcal L_a={}&F_a+\bm z_t^{\T}\bm m
 +\bm a^{\T}(\mathsf A-E\mathbf1)\bm r
 +\bm b^{\T}(\mathsf A^{\T}-E\mathbf1)\bm l
 +\nu(\bm l^{\T}\bm r-1).
 \label{eq:excfunctional}
\end{align}
Its stationarity conditions are
\begin{align}
 D_q\left[F_{a,r}+(\mathsf A-E\mathbf1)^{\T}\bm a+\nu\bm l\right]&=0,
                                                        \label{eq:adjr}\\
 F_{a,l}+(\mathsf A-E\mathbf1)\bm b+\nu\bm r&=0,
                                                        \label{eq:adjl}\\
 F_{a,E}-\bm a^{\T}\bm r-\bm b^{\T}\bm l&=0,
                                                        \label{eq:adjenergy}\\
 (J^{\T}\bm z_t)_I+(F_{a,t})_I+
   \bm a^{\T}\mathsf A^I\bm r+
   \bm l^{\T}\mathsf A^I\bm b&=0.
                                                        \label{eq:adjt}
\end{align}
First solve Eqs.~\eqref{eq:adjr}--\eqref{eq:adjenergy} for the state
multipliers; then solve Eq.~\eqref{eq:adjt} for cluster response. The
right CC equations do not depend on $r,l,E$, giving this triangular order.
The diagonal/bordered matrices used in a GMRES preconditioner are not
substitutes for the exact operators in these equations.

Multiplying Eq.~\eqref{eq:adjl} by $\bm l^{\T}$ gives useful identities:
\begin{equation}
 \nu_b=-E_\mu,\qquad \nu_c=-\delta_\mu.
 \label{eq:nuidentities}
\end{equation}
For the base objective, $\bm a=\bm b=0$, $\nu_b=-E_\mu$ solves the state
part on shell. Thus the base cluster multiplier reduces to
\begin{equation}
 J^{\T}\bm z_{t,b}=-\bm q_t,\qquad
 (q_t)_I=\me{L_\mu}{H^I}{R_\mu}.
 \label{eq:baseeomsimplified}
\end{equation}
This is the familiar fixed-T biorthogonal eigenvalue derivative supplemented
by ground-cluster response. The code's general constrained solve reproduces
this reduction. For the Q correction the state multipliers are generally
nonzero; the projected left EOM vector alone is not its gradient adjoint.

\subsection{Excited Hbar coefficients and a denominator sign table}
Embed $\bm r,\bm l,\bm a,\bm b$ in the full determinant support, and embed
$\bm z_t$ on P. The objective coefficients are
\begin{align}
 \mathcal A_b^{\rm raw}&=\bm l\bm r^{\T},\\
 \mathcal A_c^{\rm raw}&=\bm v_Q\bm r^{\T}+\bm l\bm u_Q^{\T}
                  +\mathcal D_\mu,\\
 \mathcal D_\mu^{\rm full}&=\sum_K w_K\ket{K}\bra{K},\qquad
 \mathcal D_\mu^{\rm mp}=W_Q\ket{0}\bra{0}.
 \label{eq:excrawcoeff}
\end{align}
For either objective add
\begin{equation}
 \boxed{\mathcal A_a=\mathcal A_a^{\rm raw}
    +\bm z_{t,a}\evec^{\T}+\bm a_a\bm r^{\T}
                                  +\bm l\bm b_a^{\T}.}
 \label{eq:exccoeff}
\end{equation}
Terms containing $E$ and $\nu$ in Eq.~\eqref{eq:excfunctional} are constant
with respect to $\bH$ at fixed response variables. Their effect is retained
through the adjoint equations, particularly Eq.~\eqref{eq:adjenergy}.

\begin{table}[htbp]\centering\small
\begin{tabular}{llll}
\toprule
Case & Extra Hbar diagonal coefficient & $\partial_E\delta$ & Explicit Fock term\\
\midrule
Ground, full & $+w_K$ on Q; $-W_Q$ on $00$ & Not independent & None\\
Ground, MP & None & Not independent & $-\sum_Kw_K\chi_{Kp}$\\
Excited, full & $+w_K$ on Q & $-W_Q$ & None\\
Excited, MP & $+W_Q$ on $00$ & $-W_Q$ & $-\sum_Kw_K\chi_{Kp}$\\
\bottomrule
\end{tabular}
\caption{Denominator contributions at fixed T/R/L and independent excited E.
Moment derivatives are additional in every row.}\label{tab:signs}
\end{table}

\section{Reduced total-gradient implementation for an excited state}
\label{sec:lean}
The separate stationary functionals are useful for derivation and diagnostic
component checks. They do not require two complete production response and
density calculations when only the total gradient is requested. We now
reduce their sum algebraically, keeping the same full moment energy,
independent-energy denominator response, and reference+$P$ eigenvectors.
This section describes a verified algebraic route; it has not replaced the
current oracle and is not a measured production speedup.

\subsection{Projected correction responses}
For a fixed isolated excited root, suppress $\mu$ on the vectors and set
\begin{equation}
 \mathsf B_\mu=\mathsf A-E_\mu\mathbf1,\qquad
 g_r=\bH_{Q,\{0,P\}}^{\T}\bm v,\qquad
 g_l=\bH_{\{0,P\},Q}\bm u.
 \label{eq:leanB}
\end{equation}
The operator $\mathsf B_\mu$ is the shifted projected EOM operator; it is
not the ground CC residual Jacobian J, and is not the ground-state B in
Eq.~\eqref{eq:JBe}. All state and response vectors in this section have
$1+|P|$ components, including the reference component.

The correction is separately linear in $\bm l$ and $\bm r$ at fixed
$T,E,f$, hence
\begin{equation}
 \bm l^{\T}g_l=\bm r^{\T}g_r=\delta_\mu.
 \label{eq:leansolvability}
\end{equation}
Define right and left projected responses by
\begin{align}
 \mathsf B_\mu\bm x&=-(g_l-\delta_\mu\bm r),
       &\bm l^{\T}\bm x&=0,\\
 \mathsf B_\mu^{\T}\bm y&=-(g_r-\delta_\mu\bm l),
       &\bm r^{\T}\bm y&=0.
 \label{eq:leanresponses}
\end{align}
Equation~\eqref{eq:leansolvability} guarantees that the first right-hand
side is orthogonal to the left null vector and the second to the right
null vector. An ordinary inverse of $\mathsf B_\mu$ is inappropriate:
its null vector is the target eigenstate. Use a projected or bordered solve.
For example, define
\begin{equation}
 \mathsf K_\mu=\begin{pmatrix}\mathsf B_\mu&\bm r\\
                                  \bm l^{\T}&0\end{pmatrix}.
\end{equation}
Then the two responses follow from the same bordered operator and its
transpose:
\begin{align}
 \mathsf K_\mu\begin{pmatrix}\bm x\\\xi\end{pmatrix}
 &=\begin{pmatrix}-(g_l-\delta_\mu\bm r)\\0\end{pmatrix},\\
 \mathsf K_\mu^{\T}\begin{pmatrix}\bm y\\\eta\end{pmatrix}
 &=\begin{pmatrix}-(g_r-\delta_\mu\bm l)\\0\end{pmatrix}.
 \label{eq:leanborders}
\end{align}
At a consistent converged solution $\xi=\eta=0$. The borders impose the
normalization-side conditions without requiring a dense projected matrix.
In a selected-tensor implementation, operator actions and preconditioning
can reuse the corresponding EOM right/left machinery.

\subsection{Recovery of the original adjoints and gauge}
For the total objective $F=F_b+F_c$, the on-shell state sources are
\begin{equation}
 F_r=E_\mu\bm l+g_r,\quad F_l=E_\mu\bm r+g_l,
 \quad F_E=-W_Q,
 \quad \nu_{\rm tot}=-E_\mu-\delta_\mu.
 \label{eq:leantotalsources}
\end{equation}
A convenient representative of the original adjoints is
\begin{equation}
 \bm a=\bm y-W_Q\bm l,\qquad \bm b=\bm x.
 \label{eq:leanadjoints}
\end{equation}
Substitution into Eqs.~\eqref{eq:adjr}--\eqref{eq:adjenergy} gives
\begin{align}
 \mathsf B_\mu^{\T}\bm a+F_r+\nu_{\rm tot}\bm l
 &=\mathsf B_\mu^{\T}\bm y+g_r-\delta_\mu\bm l=0,\\
 \mathsf B_\mu\bm b+F_l+\nu_{\rm tot}\bm r
 &=\mathsf B_\mu\bm x+g_l-\delta_\mu\bm r=0,\\
 \bm a^{\T}\bm r+\bm b^{\T}\bm l&=-W_Q.
 \label{eq:leanadjointproof}
\end{align}
The last equality uses $\bm r^{\T}\bm y=\bm l^{\T}\bm x=0$ and
$\bm l^{\T}\bm r=1$. It preserves the denominator's independent-energy
response exactly; that response has not been frozen or discarded.

The representative in Eq.~\eqref{eq:leanadjoints} need not have $b_q=0$.
To recover the original pivot convention, set
\begin{equation}
 \theta=\frac{x_q}{r_q},\qquad
 \bm a'=\bm a+\theta\bm l,\qquad
 \bm b'=\bm b-\theta\bm r.
 \label{eq:leanpivot}
\end{equation}
Then $b'_q=0$, the energy equation is unchanged, and
\begin{equation}
 \bm a'\bm r^{\T}+\bm l(\bm b')^{\T}
 =\bm a\bm r^{\T}+\bm l\bm b^{\T}.
\end{equation}
The corresponding cluster source is invariant for the same reason. Thus
the projected-response gauge and the original constrained pivot gauge give
the identical stationary Hbar coefficient and derivative.

For the base objective alone, Eq.~\eqref{eq:baseeomsimplified} already
provides $\bm a_b=\bm b_b=0$, $\nu_b=-E_\mu$. There is no need to
numerically solve a generic base-state bordered response problem.

\subsection{One combined cluster adjoint and three density factors}
Embed all state/response vectors on reference+$P$ and the divided moments
on Q. Add the base and correction coefficients in
Eqs.~\eqref{eq:excrawcoeff} and \eqref{eq:exccoeff}, then use
Eq.~\eqref{eq:leanadjoints}. Before cluster response, the coefficient is
\begin{align}
 \mathcal A_{\rm pre}
 &=\bm l\bm r^{\T}+\bm v_Q\bm r^{\T}+\bm l\bm u_Q^{\T}
   +(\bm y-W_Q\bm l)\bm r^{\T}+\bm l\bm x^{\T}
   +\mathcal D_\mu\\
 &=\big[(1-W_Q)\bm l+\bm y+\bm v_Q\big]\bm r^{\T}
                 +\bm l(\bm x+\bm u_Q)^{\T}+\mathcal D_\mu.
 \label{eq:leanpre}
\end{align}
Here
\begin{equation}
 \mathcal D_\mu^{\rm full}=\sum_{K\in Q}w_K\ket K\bra K,
 \qquad \mathcal D_\mu^{\rm mp}=W_Q\ket0\bra0.
 \label{eq:leandiagonal}
\end{equation}
The factor $1-W_Q$ is the chain-rule contribution from the uncorrected
state energy in the denominator. It is not an additional energy
renormalization and does not change Eq.~\eqref{eq:energy}.

A single combined cluster multiplier now solves
\begin{equation}
 \boxed{J^{\T}\bm z_t=-\bm q_t,\qquad
 (q_t)_I=\sum_{AB}(\mathcal A_{\rm pre})_{AB}H^I_{AB}.}
 \label{eq:leancluster}
\end{equation}
All state vectors, divided moments, weights and adjoints in this contraction
are held fixed. Their implicit response has already been eliminated by
stationarity. The final coefficient is
\begin{equation}
 \boxed{\mathcal A_{\rm tot}
 =\big[(1-W_Q)\bm l+\bm y+\bm v_Q\big]\bm r^{\T}
 +\bm l(\bm x+\bm u_Q)^{\T}+\bm z_t\evec^{\T}+\mathcal D_\mu.}
 \label{eq:leanthreefactors}
\end{equation}
There are three outer products plus the separate denominator diagonal.
Their pullback and transition-density contractions are those of
Sec.~\ref{sec:densities}. Multiplying two effective ``corrected'' state
vectors together would introduce cross terms absent from
Eq.~\eqref{eq:leanthreefactors}; that is not this reduction.
For MP, the explicit $d_p=-\sum_Kw_K\chi_{Kp}$ Fock density in
Sec.~\ref{sec:fock} must still be added.

\subsection{Production work sequence and reuse}
For a total gradient, the resulting sequence at a fixed geometry is:
\begin{enumerate}
 \item Reuse converged T, Hbar and the target right/left eigenstate from the
       uncorrected energy calculation. Preserve the reference component.
 \item Construct Q moments, denominators and response sources in bulk.
       No response solve is required per Q determinant or Cartesian coordinate.
 \item Solve the two projected state equations in Eq.~\eqref{eq:leanresponses}.
 \item Form the merged source and solve the single cluster equation in
       Eq.~\eqref{eq:leancluster}.
 \item Accumulate one total density from the three factors and denominator
       terms, then solve one total orbital response and contract one set of
       AO derivatives, including overlap and fitted-metric terms.
\end{enumerate}
For fixed SD/selected-T ranks, a production realization should evaluate these
objects through connected tensor or selected-list kernels. Full-CI support
and dense full-space coefficient matrices are not required by the equations.
An exact asymptotic tensor scaling or wall-time speedup requires a separate
contraction/implementation audit.

At a fixed geometry and selection, T, Hbar, J and the HF response operator
can be shared across roots and denominator variants. R/L are shared across
denominator variants for a given root. Different denominators change the
correction sources, so their response solutions are not interchangeable.
The shifted/projected EOM operator also remains root-specific, although its
underlying action and preconditioning structure can be reused.

If separate gradient components are needed, compute the base response once
and cache it. Its state adjoint is analytic; only its cluster response and
density need to be evaluated. Otherwise, the combined route avoids separate
base/correction cluster solves, density builds and orbital-response solves.
This changes the implementation schedule, not the two-functional theory.

\subsection{Algebraic verification and current implementation status}
\label{sec:leanvalidation}
The prototype compares the merged total bare-Hamiltonian gradient
coefficient with the independent
dense constrained oracle for closed/nonclosed P, two excited roots, and
full/MP denominators: eight cases. The maximum coefficient discrepancy is
$\ee{5.551}{-15}$. The checker and saved results are:
\begin{quote}\small
\path{devtools/check_ccpq_lean_response.py}\\
\path{docs/theory/ccpq_lean_response_checks.json}.
\end{quote}
The explicit MP Fock density is retained separately; coefficient agreement
at fixed Fock must not be misreported as permission to omit it.

The original molecular tables continue to document the existing oracle's
validated gradients. They are not timings or independent molecular tests
of a newly optimized implementation. The reduction above is algebraically
verified, but has not replaced the production/oracle path; no measured
speedup or globally minimal implementation is claimed. A production change
should retain the original route as a correctness reference, compare the
resulting densities and nuclear gradients, and record phase-resolved timings.


\section{Pullback to bare-integral derivative densities}
\label{sec:densities}
Use the explicit pairing
$\langle X,Y\rangle_F=\sum_{AB}X_{AB}Y_{AB}$. This avoids an implicit
transpose in density definitions. At fixed multipliers,
\begin{equation}
 \dd\mathcal L_a\big|_{T,r,l,\ldots}
  =\langle\mathcal A_a,U\dd H V\rangle_F+
                                     \dd\mathcal L_a\big|_{f},
\end{equation}
and rearranging indices gives
\begin{equation}
 \boxed{\mathcal K_a=U^{\T}\mathcal A_a V^{\T},\qquad
       \dd\mathcal L_a\big|_H=\sum_{AB}(\mathcal K_a)_{AB}\dd H_{AB}.}
 \label{eq:pullback}
\end{equation}
For an outer product $\mathcal A=\bm u\bm v^{\T}$,
$\mathcal K=(U^{\T}\bm u)(V\bm v)^{\T}$: a transition density between two
vectors suffices. Diagonal Q weights are streamed or handled by exact sparse
commutator expressions; no dense full-space coefficient matrix is required.

In spin orbitals, take
\begin{equation}
 H=\sum_{pq}h_{pq}a_p^\dagger a_q+
   \frac14\sum_{pqrs}\langle pq\Vert rs\rangle
                      a_p^\dagger a_q^\dagger a_s a_r.
 \label{eq:Hsecond}
\end{equation}
The derivative densities are
\begin{align}
 \gamma_{pq}^{a}&=\sum_{AB}(\mathcal K_a)_{AB}
                         \me{A}{a_p^\dagger a_q}{B},\\
 \Gamma_{pqrs}^{a}&=\sum_{AB}(\mathcal K_a)_{AB}
                         \me{A}{a_p^\dagger a_q^\dagger a_s a_r}{B},\\
 \dd\mathcal L_a\big|_H&=\sum_{pq}\gamma_{pq}^{a}\dd h_{pq}
       +\frac14\sum_{pqrs}\Gamma_{pqrs}^{a}\dd\langle pq\Vert rs\rangle.
 \label{eq:spinrdms}
\end{align}
These are derivative densities of a stationary functional, not necessarily
Hermitian expectation-value densities of a variational wavefunction. Only
the symmetry components seen by the real integrals contribute.

\subsection{The explicit MP denominator density}
\label{sec:fock}
Define
\begin{equation}
 d_p=-\sum_{K\in Q}w_K\chi_{Kp},\qquad
 f_{pp}=h_{pp}+\sum_i n_i^0\langle pi\Vert pi\rangle.
 \label{eq:dfock}
\end{equation}
Then
\begin{equation}
 \dd\mathcal L_c\big|_f=\sum_p d_p\dd f_{pp}
 =\sum_p d_p\dd h_{pp}+\sum_{pi}d_pn_i^0\dd\langle pi\Vert pi\rangle.
 \label{eq:fockchain}
\end{equation}
Let $d_{pq}=d_p\delta_{pq}$ and $n_{pq}=n_p^0\delta_{pq}$. The required
additions to Eq.~\eqref{eq:spinrdms} are
\begin{align}
 \gamma^D_{pq}&=d_{pq},\\
 \Gamma^D_{pqrs}&=d_{pr}n_{qs}-d_{ps}n_{qr}
                         -d_{qr}n_{ps}+d_{qs}n_{pr}.
 \label{eq:denrdm}
\end{align}
Each antisymmetric image contributes the same contraction, so the
$\tfrac14\Gamma^D:\dd v$ expression reduces to the second term of
Eq.~\eqref{eq:fockchain}. This is why a Hamiltonian coefficient evaluated
at fixed $f$ is insufficient for an MP gradient.

Since every K has the same electron number as the reference,
$\sum_p\chi_{Kp}=0$, and hence $\sum_p d_p=0$. A uniform shift of all Fock
diagonal entries therefore leaves the denominator unchanged. Although a
frozen-core orbital has $\chi_{Kp}=0$, its occupation still enters $n_i^0$
in Eq.~\eqref{eq:fockchain}; deleting core spectator terms is incorrect.

\subsection{Spatial convention and spin tracing}
For a shared spatial basis,
\begin{align}
 D_{pq}&=\sum_\sigma\gamma_{p\sigma,q\sigma},\\
 G_{pqrs}&=\sum_{\sigma\tau}
                  \Gamma_{p\sigma,q\tau,r\sigma,s\tau},\\
 \dd\mathcal L&=\sum_{pq}D_{pq}\dd h_{pq}
                       +\frac12\sum_{pqrs}G_{pqrs}\dd(pr|qs).
 \label{eq:spatial}
\end{align}
Here $(ab|cd)$ is a chemists' integral. The ordering $(pr|qs)$ in
Eq.~\eqref{eq:spatial} is deliberate. Confusing it with $(pq|rs)$ without
transposing G is a two-particle-density convention error. The code averages
D and G over the symmetries actually seen by these real integrals before
orbital response.

\section{Canonical RHF orbital relaxation and the overlap term}
\label{sec:orbital}
The amplitudes have been relaxed, but their densities are not stationary
with respect to the reference orbitals. This section derives the remaining
response for the restricted, common-orbital reference. The same linear
operations apply to the base and correction densities separately.

Let $C$ denote the AO-to-MO coefficient matrix and $P_o$ the \emph{per-spin}
occupied MO projector (one on each occupied spatial orbital, not two).
Define the RHF response map
\begin{equation}
 \mathcal R(X)_{pq}=\sum_{rs}[2(pq|rs)-(pr|qs)]X_{rs},\qquad
 f=h+\mathcal R(P_o).
 \label{eq:rhfresponse}
\end{equation}
The generalized Fock matrix of the amplitude-relaxed functional is
\begin{equation}
 \mathcal F_{tp}=\sum_qh_{tq}D_{pq}
                 +\sum_{qrs}(tr|qs)G_{pqrs},\qquad
 g_{tp}=2(\mathcal F_{tp}-\mathcal F_{pt}).
 \label{eq:orbitalgradient}
\end{equation}
Indeed, under $C\mapsto C(\mathbf1+\kappa)$ with $\kappa^{\T}=-\kappa$,
\begin{equation}
 \dd h=\kappa^{\T}h+h\kappa,\qquad
 \dd\mathcal L=2\sum_{tp}\mathcal F_{tp}\kappa_{tp}
              =\sum_{t>p}g_{tp}\kappa_{tp}.
\end{equation}
The analogous four-index integral rotation supplies the two-electron part
of $\mathcal F$ in Eq.~\eqref{eq:orbitalgradient}.

\subsection{Canonical constraints and the Z vector}
The reference satisfies the occupied--virtual Brillouin conditions and
canonical off-diagonal Fock conditions. Selected spaces and orbital-energy
denominators can make rotations within occupied or virtual subsets matter.
Frozen-core/active-occupied rotations also matter because core excitation
amplitudes are excluded. Introduce a multiplier $\zeta_{pq}$ for each
independent relevant $f_{pq}=0$ and collect it in a symmetric matrix S with
$S_{pq}=S_{qp}=\zeta_{pq}$, zero diagonal. Write $Z=S/2$ so the added term is
\begin{equation}
 \mathcal L_{\rm orb}=\mathcal L+\sum_{p>q}\zeta_{pq}f_{pq}
                    =\mathcal L+\sum_{pq}Z_{pq}f_{pq}.
 \label{eq:orblagrangian}
\end{equation}
For rotations within the same occupied or virtual block, the reference
projector does not change. Therefore
\begin{equation}
 \zeta_{pq}=-\frac{g_{pq}}{\epsilon_p-\epsilon_q}
 \label{eq:zetainternal}
\end{equation}
for active-occupied pairs, virtual pairs and core/active-occupied pairs.
A degenerate pair with a nonzero source does not admit this division; a
selected model must then specify a differentiable orbital convention.

For a virtual $a$ and occupied $i$, the projector changes by the symmetric
OV rotation. Differentiating the Fock operator gives
\begin{equation}
 \delta f_{ai}=(\epsilon_a-\epsilon_i)\kappa_{ai}
                              +\mathcal R(S_{\rm OV}(\kappa))_{ai}.
\end{equation}
Using self-adjointness of $\mathcal R$ in Eq.~\eqref{eq:orblagrangian},
stationarity yields
\begin{equation}
 \boxed{(\epsilon_a-\epsilon_i)\zeta_{ai}
  +\mathcal R(S_{\rm OV}(\zeta))_{ai}
       =-g_{ai}-\mathcal R(S_{\rm int})_{ai}.}
 \label{eq:orbz}
\end{equation}
$S_{\rm int}$ contains the multipliers in Eq.~\eqref{eq:zetainternal}.
All occupied orbitals, including core, participate in the OV response.
Rotations purely within a completely frozen core have zero source for the
stated model. Core/active response must not be discarded merely because
core excitation amplitudes were frozen.

\subsection{Relaxed densities}
At fixed $Z$, differentiating the extra Fock term with respect to bare
integrals gives
\begin{align}
 D^{\rm rel}_{pq}&=D_{pq}+Z_{pq},\\
 G^{\rm rel}_{pqrs}&=G_{pqrs}
   +2(Z_{pr}(P_o)_{qs}+(P_o)_{pr}Z_{qs})
   -(Z_{ps}(P_o)_{qr}+(P_o)_{ps}Z_{qr}).
 \label{eq:relaxeddensities}
\end{align}
Substitution in $D^{\rm rel}:\dd h+\tfrac12G^{\rm rel}:\dd g$
produces $D:\dd h+\tfrac12G:\dd g+Z:\dd f$ exactly. Thus the orbital
multipliers are incorporated into the same derivative-density convention.

\subsection{AO overlap and the energy-weighted density}
Differentiate $C^{\T}S_{\rm AO}C=\mathbf1$. Writing
$C^x=C\bigl(\kappa^x-\tfrac12S^x_{\rm MO}\bigr)$ separates the antisymmetric
orbital rotation from the symmetric orthonormality connection, with
$S^x_{\rm MO}=C^{\T}S^x_{\rm AO}C$. The antisymmetric terms cancel by
Eq.~\eqref{eq:orbz}. The symmetric connection contributes
$-\sym(\mathcal F):S^x_{\rm MO}$ from the amplitude functional.
The explicit Fock rotation contributes
$-\tfrac12(\epsilon Z+Z\epsilon):S^x_{\rm MO}$.
Finally the occupied projector changes by
$-\tfrac12(S^x_{\rm MO}P_o+P_oS^x_{\rm MO})$, giving
$-\tfrac12[P_o\mathcal R(Z)+\mathcal R(Z)P_o]:S^x_{\rm MO}$.
Hence
\begin{equation}
 W=\sym(\mathcal F)+\frac12(\epsilon Z+Z\epsilon)
            +\frac12[P_o\mathcal R(Z)+\mathcal R(Z)P_o],
 \qquad \epsilon=\diag(\epsilon_p).
 \label{eq:weighted}
\end{equation}
This overlap coefficient is not supplied by the amplitude response alone.

\section{The final nuclear gradients}
\subsection{Conventional four-center integrals}
Back-transform every spatial index with C, for example
$D^{\rm AO}=CD^{\rm rel}C^{\T}$ and $W^{\rm AO}=CWC^{\T}$. Then
\begin{equation}
 \boxed{\mathcal E_\mu^x=V_{\rm nn}^x+
  \sum_{ab}D^{\rm AO}_{ab}h^x_{ab}
  -\sum_{ab}W^{\rm AO}_{ab}S^x_{ab}
  +\frac12\sum_{abcd}G^{\rm AO}_{abcd}(ac|bd)^x.}
 \label{eq:final4c}
\end{equation}
All AO derivatives here are evaluated at fixed MO coefficients. Motion of
the orbitals has already been included by relaxation and the overlap term.
The correction gradient uses the same contraction with correction densities
and \emph{without} $V_{\rm nn}^x$. Linearity of orbital response gives
$g_{\rm total}=g_b+g_c$; the driver evaluates total and correction, then
obtains the base by subtraction.

For corrected excitation energies one may form
\begin{equation}
 \Omega_\mu^{\rm corr}=\mathcal E_\mu-\mathcal E_0,
 \qquad (\Omega_\mu^{\rm corr})^x=g_\mu^x-g_0^x.
\end{equation}
The molecular tables in this note compare total \emph{state}-energy gradients;
they should not be relabeled excitation-energy-gradient tests.

\subsection{Density fitting and Cholesky representations}
Let the fitted spatial integral be
\begin{equation}
 (pr|qs)_{\rm fit}=\sum_X B_{Xpr}B_{Xqs}.
\end{equation}
For a G symmetrized over the real integral pair symmetries,
\begin{equation}
 \frac12\sum_{pqrs}G_{pqrs}\dd(pr|qs)_{\rm fit}
   =\sum_{Xpr}Y_{Xpr}\dd B_{Xpr},\qquad
 Y_{Xpr}=\sum_{qs}G_{pqrs}B_{Xqs}.
 \label{eq:Ydensity}
\end{equation}
For the common restricted orbital variation, the native per-spin machinery
can use $D^a=D^b=D/2$ and $Y^a=Y^b=Y/2$. This is a bookkeeping split of the
shared-orbital derivative, not a reconstruction of arbitrary spin-resolved
RDMs. The native orbital response then returns relaxed $D^s,Y^s,W^s$.

To expose the auxiliary-basis derivative, let
\begin{equation}
 I_{Qab}=(Q|ab),\quad M_{QR}=(Q|R),\quad M=LL^{\T},\quad
 X=L^{-1},\quad B_{Xab}=\sum_Q X_{XQ}I_{Qab}.
\end{equation}
After AO back-transformation and summation over spins, define
\begin{align}
 (W_3)_{Qab}&=\sum_X X_{XQ}Y^{\rm AO}_{Xab},\\
 C_{XY}&=\sum_{s,pq}B^s_{Xpq}Y^s_{Ypq},\qquad
 W_2=-\frac12X^{\T}CX.
 \label{eq:dfweights}
\end{align}
The metric term follows from $XMX^{\T}=\mathbf1$: the symmetric part of
$\dd X\,X^{-1}$ is $-\tfrac12X(\dd M)X^{\T}$. Auxiliary orthogonal
rotations do not change the fitted ERI functional, so only that symmetric
part survives its contraction with the symmetric part of C. Thus
\begin{equation}
 \boxed{\mathcal E_\mu^x=V_{\rm nn}^x+D^{\rm AO}:h^x-W^{\rm AO}:S^x
                            +W_3:I^x+W_2:M^x.}
 \label{eq:finaldf}
\end{equation}
Both three-center and metric derivatives include auxiliary-center motion.
Using a DF energy but the conventional four-center derivative in
Eq.~\eqref{eq:final4c} would compare two Hamiltonians. In CCpy,
\texttt{integral\_model="factorized"} differentiates the fitted model;
\texttt{"conventional"} differentiates the unfitted model.
For the Cholesky derivative, Q labels locally fixed pivot pairs and the same
metric construction applies. Discrete pivot reselection is not itself
analytically differentiated.

\section{Determinant implementation and equation-to-code correspondence}
\label{sec:implementation}
The independent dense oracle uses complex-step differentiation of the
constraint system and solves its transposed Jacobian. The larger backend
uses Slater-rule Hamiltonian actions from the external FCIpy routines,
Davidson right/left roots, and GMRES response. It does not receive optimized
CCpy amplitudes or eigenvectors. With connected sparse caching it stores a
sparse bare Hamiltonian, not a dense determinant-pair Hamiltonian.

For a factor pair $\bm u\bm v^{\T}$, the program evaluates the T source as
\begin{equation}
 \partial_{t_I}(\bm u^{\T}\bH\bm v)
   =(\bH^{\T}\bm u)^{\T}\tau_I\bm v
                         -\bm u^{\T}\tau_I(\bH\bm v).
 \label{eq:factorgradient}
\end{equation}
Sparse excitation transitions contract this expression for all I without
forming a dense third-order derivative tensor. Its bare-integral density
is a transition RDM of $U^{\T}\bm u$ and $V\bm v$, as in
Eq.~\eqref{eq:pullback}. Exact T/R/L response is still solved; using a sparse
representation does not introduce a two-body approximation to the moments.

\begin{longtable}{>{\raggedright\arraybackslash}p{0.28\textwidth}>{\raggedright\arraybackslash}p{0.65\textwidth}}
\toprule Construction & Implementation location and role\\\midrule\endhead
Definitions and phases & \path{ccpy/gradients/ccpq/determinants.py},
\path{spaces.py}: normalized determinant excitations, fixed P/Q selections,
CCpy triples-list conversion.\\
Ground functionals & \path{amplitudes.py}, \texttt{SelectedCC.run}:
$J,B,\lambda^g,\lambda,\sigma,\pi$ and Eq.~\eqref{eq:groundcoeff}.\\
Unified dense adjoint & \path{state_oracle.py}, \texttt{evaluate/solve}:
Eqs.~\eqref{eq:masterfunctional}, \eqref{eq:excconstraints} and
\eqref{eq:objectives}; independent complex-step Jacobians and coefficients.\\
Matrix-free adjoint & \path{matrix_free.py}, \texttt{objective},
\texttt{constraint\_factors}, \texttt{vjp}, \texttt{solve\_adjoint}:
Eqs.~\eqref{eq:excsources}--\eqref{eq:exccoeff}, including the sign table.\\
Densities & \path{densities.py}: Eqs.~\eqref{eq:pullback},
\eqref{eq:spinrdms}, \eqref{eq:denrdm} and \eqref{eq:spatial}.\\
Hamiltonian and RDM actions & \path{fcipy_adapter.py},
\path{fortran/fcipy_connected.f90}: connected Slater rules, transition RDMs,
and exact refresh of numerical Hamiltonian entries at each geometry.\\
Conventional orbitals/AO derivative & Local \path{orbital_response.py},
\path{nuclear_gradient.py}: Eqs.~\eqref{eq:orbz}--\eqref{eq:final4c}.\\
DF orbitals/AO derivative & Parent \path{ccpy/gradients/orbital_response.py}
and \path{derivative_integrals.py}, used through \texttt{driver.py}:
Eq.~\eqref{eq:finaldf}.\\
Independent optimized comparator & \texttt{CRCC23},
\texttt{CCP3(two\_body\_approx=False)}, \texttt{EOMCCSDT\_P} and
\texttt{EOMCCP3}, with the identical determinant lists and integral model.\\
\bottomrule
\end{longtable}
Paths without the full prefix in the first seven rows are relative to
\path{ccpy/gradients/ccpq/}; the last orbital row explicitly refers to the
parent gradient package.

\subsection{The ground and excited approximate APIs are different models}
Identical P/Q sets do not specify all moment approximations. The default
optimized \texttt{CCT3} uses CCSD-like correction contractions and left-CCSD.
The matching full ground specialization for a CCSDt P list is
\texttt{CCP3(two\_body\_approx=False)}. Their energy and gradient differences
were measured, not declared a passing comparison of identical functionals.

In the excited code, \texttt{EOMCCT3(two\_body\_approx=False)} retains VT3
contributions in Hbar, but its correction kernels still use only the SD
parts of R and L. Full selected-R3/L3 moments are supplied by
\texttt{EOMCCP3}. That is the full optimized excited-state comparator in the
tables below. Agreement with it does not establish an analytical gradient
for the separate SD-only correction approximation.

\subsection{Practical algorithm}
\begin{enumerate}
\item Fix reference, core convention, P/Q lists, denominator model and target root.
\item Solve canonical RHF in the same conventional or fitted integral model
      that will be differentiated; transform the bare Hamiltonian.
\item Solve right CC(P). Construct the projected ground left state and
      $\lambda^g$, or the normalized projected excited right/left root.
\item Form Q moments, denominators, $u,v,w$, and the correction energy.
\item For component diagnostics, solve the base and correction adjoints
      separately. For an excited-state total gradient, use the projected
      responses and single combined cluster adjoint in Sec.~\ref{sec:lean}.
      Check the right/left constraints and stationarity residuals.
\item Form Hbar covectors, pull them back through $U,V$, and evaluate
      transition derivative densities. Add Eq.~\eqref{eq:denrdm} for MP.
\item Spin trace, symmetrize and solve orbital response, including canonical
      internal and frozen-core/active rotations. Construct W.
\item Contract the matching AO integral derivatives. Add nuclear repulsion
      exactly once, to the base/total rather than correction gradient.
\item For validation, re-solve the optimized energy at all displaced points,
      hold lists fixed, verify root identities and test step convergence.
\end{enumerate}

\clearpage
\section{Numerical validation and its evidentiary scope}
\label{sec:validation}
The tables below are generated from saved JSON reports rather than copied
from rounded conversational summaries. Units are $\gunit$ for errors and
$\eunit$ for energies. Every error is a maximum absolute deviation over the
\emph{explicitly stated} finite-differenced components. The independent
optimized tensor comparisons validate both amplitude/left-state equations
and derivatives; determinant-energy finite differences for more general
T/Q models separately validate differentiation of those defined models.

\subsection{Independent optimized ground-state tensor comparisons}
The primary DF water calculation freezes O(1s), retaining all 13 spatial orbitals.
CR-CC(2,3) has P = 64 S + 1,360 D and Q = all 11,200 T. The active-SDt
case has P = 64 S + 1,360 D + 2,436 T and Q = all 8,764 remaining T.
All nine components use four-point differences at 0.001 bohr of independently
converged optimized tensor energies, with the same DF auxiliary basis and reference.
The full CC(P;3) comparator includes selected triples in the left state and moments.

\begin{table}[htbp]\centering\small
\begin{tabular}{lllll}
\toprule
Method & Variant & Base & Correction & Total\\
\midrule
CR-CC(2,3) & A & $\ee{1.091}{-10}$ & $\ee{5.870}{-14}$ & $\ee{1.090}{-10}$\\
CR-CC(2,3) & D & $\ee{1.091}{-10}$ & $\ee{5.509}{-14}$ & $\ee{1.090}{-10}$\\
Full CC(P;3) & A & $\ee{1.152}{-10}$ & $\ee{3.064}{-14}$ & $\ee{1.152}{-10}$\\
Full CC(P;3) & D & $\ee{1.152}{-10}$ & $\ee{5.232}{-14}$ & $\ee{1.152}{-10}$\\
\bottomrule
\end{tabular}
\caption{DF H2O/6-31G: independent optimized-energy finite differences; maximum errors over nine components.}\label{tab:groundnumerical}
\end{table}

The determinant and existing DF CR-CC(2,3)-A analytical gradients differ by at most $\ee{1.432}{-13}$ $\gunit$.
The largest base/correction energy differences over this table are $\ee{2.842}{-13}$ / $\ee{2.489}{-15}$ $\eunit$.

\subsection{Independent optimized excited-state tensor comparisons}
The water EOM test uses the same 2,436 active triples in ground T and excited R/L,
all S/D, and all 8,764 remaining triples in Q. Both multiplicities use full spin blocks.
The lowest triplet and singlet map to roots 1 and 2 of the unfiltered fixed-$M_s=0$
spectrum. All nine components use a 0.001-bohr four-point stencil.

\begin{table}[htbp]\centering\small
\begin{tabular}{lllll}
\toprule
State & Variant & Base & Correction & Total\\
\midrule
Triplet & A & $\ee{1.220}{-10}$ & $\ee{1.504}{-13}$ & $\ee{1.220}{-10}$\\
Triplet & D & $\ee{1.220}{-10}$ & $\ee{1.777}{-12}$ & $\ee{1.202}{-10}$\\
Singlet & A & $\ee{1.271}{-10}$ & $\ee{1.992}{-13}$ & $\ee{1.271}{-10}$\\
Singlet & D & $\ee{1.270}{-10}$ & $\ee{1.580}{-12}$ & $\ee{1.254}{-10}$\\
\bottomrule
\end{tabular}
\caption{Full EOMCCSDt/EOMCC(P;3), frozen-O(1s) H2O/6-31G. The base column separately checks EOMCCSDt derivatives.}\label{tab:waterexcited}
\end{table}

Independent Arnoldi on optimized HR and determinant Davidson spectra agree within $\ee{7.816}{-13}$ $\eunit$.
Physical right/left determinant-coordinate overlaps differ from unity by at most $\ee{1.332}{-15}$.
The maximum base/correction energy differences are $\ee{1.705}{-13}$ / $\ee{1.269}{-14}$ $\eunit$.
Displaced root overlaps exceed 0.999997; energy-gap and left/right residual guards are checked at every point.

\subsection{A nonzero-reference-component excited-state control}
The nonplanar H4/6-31G test retains all four electrons and eight orbitals.
P has 24 S + 174 D + 160 T; Q has all 200 remaining T. Its lowest triplet
and lowest singlet are roots 1 and 3 of the unfiltered $M_s=0$ spectrum.
The singlet has $|r_0|\simeq0.083$. Both corrections are nonzero.
All twelve Cartesian components use a 0.00025-bohr four-point step.

\begin{table}[htbp]\centering\small
\begin{tabular}{lllll}
\toprule
State & Variant & Base & Correction & Total\\
\midrule
Triplet & A & $\ee{1.737}{-11}$ & $\ee{3.067}{-14}$ & $\ee{1.738}{-11}$\\
Triplet & D & $\ee{1.737}{-11}$ & $\ee{4.281}{-14}$ & $\ee{1.739}{-11}$\\
Singlet & A & $\ee{2.685}{-10}$ & $\ee{6.901}{-11}$ & $\ee{1.995}{-10}$\\
Singlet & D & $\ee{2.684}{-10}$ & $\ee{9.716}{-11}$ & $\ee{1.713}{-10}$\\
\bottomrule
\end{tabular}
\caption{Independent full excited-state H4/6-31G comparison, including a nonzero right reference coefficient.}\label{tab:h4excited}
\end{table}

The singlet has a nearby root separated by about 0.00228 hartree. The initial
0.001-bohr step has visible truncation error: halving it reduces the correction
error by approximately sixteen. Smaller-step agreement is obtained without
altering the analytical equations or loosening the comparison tolerance.

\begin{table}[htbp]\centering\small
\begin{tabular}{llll}
\toprule
Step (bohr) & Base & Correction & Total\\
\midrule
0.00100 & $\ee{7.155}{-8}$ & $\ee{2.486}{-8}$ & $\ee{4.669}{-8}$\\
0.00050 & $\ee{4.473}{-9}$ & $\ee{1.559}{-9}$ & $\ee{2.915}{-9}$\\
0.00025 & $\ee{2.684}{-10}$ & $\ee{9.716}{-11}$ & $\ee{1.713}{-10}$\\
\bottomrule
\end{tabular}
\caption{H4 singlet D, atom 0/z: explicit four-point step convergence; errors in $\gunit$.}\label{tab:step}
\end{table}

\subsection{Substantial triples/quadruples selections}
The all-electron H6/6-31G models retain six electrons and twelve spatial orbitals
(48,400 determinants). Every P includes all 54 S and 945 D plus the selected
T/Q below. Q is the complete remaining triples/quadruples complement.

\begin{table}[htbp]\centering\small
\begin{tabular}{lllll}
\toprule
Active virtuals & P triples & P quadruples & Q triples & Q quadruples\\
\midrule
6 & 1660 & 2745 & 4340 & 13455\\
7 & 2716 & 5439 & 3284 & 10761\\
\bottomrule
\end{tabular}
\caption{H6 complete active selections: all three occupied MOs and either six or seven virtual MOs define P-space triples and quadruples.}\label{tab:h6spaces}
\end{table}

The virtual-6 model retains 27.67\% of triples and 16.94\% of quadruples;
the virtual-7 model retains 45.27\% and 33.57\%, respectively. Only two
components (atom 1/z and atom 2/x) are directly finite-differenced, not all eighteen.
Here the energies being finite-differenced come from the determinant model itself:
this is a general-space derivative-consistency test, not an optimized quadruples-kernel comparison.

\begin{table}[htbp]\centering\small
\begin{tabular}{lllll}
\toprule
Active virtuals & State & Base & Correction & Total\\
\midrule
6 & Ground & $\ee{8.847}{-12}$ & $\ee{4.812}{-13}$ & $\ee{8.811}{-12}$\\
6 & Root 1 & $\ee{5.745}{-12}$ & $\ee{4.548}{-12}$ & $\ee{9.855}{-12}$\\
6 & Root 2 & $\ee{9.871}{-12}$ & $\ee{1.443}{-11}$ & $\ee{2.404}{-11}$\\
7 & Ground & $\ee{2.005}{-11}$ & $\ee{2.564}{-13}$ & $\ee{2.027}{-11}$\\
7 & Root 1 & $\ee{2.388}{-12}$ & $\ee{6.128}{-12}$ & $\ee{4.930}{-12}$\\
7 & Root 2 & $\ee{1.314}{-11}$ & $\ee{6.665}{-12}$ & $\ee{1.959}{-11}$\\
\bottomrule
\end{tabular}
\caption{H6 substantial-T/Q derivative consistency. Each row takes the worst error across MP/full over the two checked components.}\label{tab:h6errors}
\end{table}

The conventional frozen-O(1s) water T/Q model has 64 S + 1,360 D +
2,480 T + 5,210 quadruples in P, and 8,720 remaining T + 37,490 remaining
quadruples in Q (245,025 support determinants). Only atom 1/z and atom 2/y
are directly finite-differenced; all nine analytical components are recorded.
No all-nine-component claim is made for this one-frozen-orbital T/Q model.

\begin{table}[htbp]\centering\small
\begin{tabular}{llll}
\toprule
State & Base & Correction & Total\\
\midrule
Ground & $\ee{4.878}{-11}$ & $\ee{1.217}{-12}$ & $\ee{4.987}{-11}$\\
Root 1 & $\ee{1.312}{-10}$ & $\ee{5.604}{-13}$ & $\ee{1.429}{-10}$\\
Root 2 & $\ee{8.063}{-11}$ & $\ee{1.995}{-11}$ & $\ee{7.836}{-11}$\\
\bottomrule
\end{tabular}
\caption{Conventional water general T/Q model: worst errors over MP/full and the two checked components.}\label{tab:watertq}
\end{table}

A separate reduced-correlation water regression freezes three occupied
orbitals (six electrons), correlating four electrons in the same 13-orbital basis.
Its 2,025-determinant model checks all nine components for ground and two excited
states with both denominators. The largest total/correction errors are $\ee{8.125}{-10}$ / $\ee{2.535}{-11}$ $\gunit$.
This is not the standard one-frozen-O(1s) model.



\subsection{What these results support}
The optimized SD/SDt cases independently reproduce the stated base and
correction energies and their state gradients. The general selected-T/Q
cases extend derivative consistency to sizeable higher-rank P and Q spaces.
The nonclosed-P and gauge audits test algebra that is not exercised simply
by choosing all S/D plus triples. Dense complex-step and matrix-free
implementations provide different differentiation paths.

This evidence does not prove the quantitative physical accuracy of a
noniterative energy correction, universal root following through crossings,
or every denominator approximation used under a CR-CC name. The molecular
orbital layer tested here is RHF. Water/6-31G results are explicitly
frozen-core, with the separate more-reduced regression identified below;
no all-electron water result is claimed. The formal model is smooth only
on a fixed-selection, nondegenerate branch with solvable response and
nonzero denominators. A numerical match is not a mathematical proof that
all future implementations or cases are error-free.

\appendix
\clearpage
\section{Exact full-denominator diagonal with high-rank cluster amplitudes}
Give a determinant its excitation rank relative to $\ket{0}$. Every pure
$T_r$ raises that rank by r when its action is nonzero. A two-body H can
lower it by at most two. Therefore a diagonal matrix element of a BCH
term can contain total excitation rank from T of at most two. Exactly,
\begin{equation}
 \bH_{KK}=H_{KK}+\me{K}{[H,T_1+T_2]}{K}
                     +\frac12\me{K}{[[H,T_1],T_1]}{K}.
 \label{eq:exactdiagonal}
\end{equation}
This rank argument is independent of whether T also contains selected
triples, quadruples or higher ranks. It explains the exact sparse diagonal
acceleration in \texttt{MatrixFreeState.prepare\_diagonal}. It is not a
truncation of the \emph{moments}: high-rank T amplitudes still affect
$M,N$, the eigenstates, and the response of $T_1,T_2$. Their total influence
on a nuclear derivative must still be retained. Tests compare this diagonal
and its amplitude/integral derivatives with dense complex-step expressions.

\section{Energy homogeneity checks and the missing-MP-term diagnostic}
Scale all bare electronic integrals by a common factor with orbitals and
occupations fixed. T and normalized isolated eigenvectors are unchanged
on the corresponding branch; E, f, D and $\delta$ scale linearly. Euler's
identity for the amplitude-relaxed densities is consequently
\begin{equation}
 \sum_{pq}\gamma^b_{pq}h_{pq}+\frac14\sum_{pqrs}\Gamma^b_{pqrs}v_{pqrs}=E_\mu,
 \quad
 \sum_{pq}\gamma^c_{pq}h_{pq}+\frac14\sum_{pqrs}\Gamma^c_{pqrs}v_{pqrs}=\delta_\mu.
 \label{eq:euler}
\end{equation}
These checks use \emph{amplitude-relaxed}, not orbital-relaxed, densities.
The latter also contain reference-orbital constraint multipliers and are
not to be interpreted as the same energy-reconstruction test.

For ground full denominators, the raw coefficient in
Eq.~\eqref{eq:groundcoeff} contracts to
\begin{equation}
 2\delta+\sum_Kw_K(\bH_{KK}-E_0)=2\delta-\delta=\delta.
\end{equation}
The Pi and Sigma constraint terms vanish on shell. For ground MP, the raw
Hbar coefficient instead gives $2\delta$, and
$\sum_p d_pf_{pp}=-\sum_Kw_KD_K=-\delta$ supplies the missing part.

For an excited full denominator,
$\langle\mathcal A_c^{\rm raw},\bH\rangle_F=\delta+EW_Q$.
For MP, adding the explicit Fock contribution to the raw Hbar contraction
gives the same $\delta+EW_Q$. Equation~\eqref{eq:adjenergy} requires
$\bm a^{\T}\bm r+\bm b^{\T}\bm l=-W_Q$, so the constraint coefficient
contracts to $-EW_Q$. This recovers Eq.~\eqref{eq:euler} in both cases.
It also pinpoints why differentiating only an Hbar coefficient at fixed
Fock diagonal fails the MP test.

\section{Gauge, empty-Q and complete-P checks}
Under $\bm r\mapsto c\bm r$, $\bm l\mapsto\bm l/c$, the moments transform
as $M\mapsto cM$, $N\mapsto N/c$, while $D,w,\delta$ are unchanged. The
full stationary derivative must be invariant too; gauge constraints change
coordinates and multipliers, not the energy surface. The implementation
tests positive and negative nonzero pivot scales.

For empty Q, $u,v,w,\delta$ and the correction sources vanish; all correction
multipliers and densities can be zero, leaving the appropriate CC(P) or
EOMCC(P) gradient. If P spans every allowed excited determinant of the
chosen frozen-core Hilbert space, Q is empty and the physical ground
branch gives the full-CI result in that support, provided the exact state
has nonzero reference overlap and a regular CC representation. No claim is made that an
arbitrary nonlinear CC initial guess always converges to that branch.

\section{Manuscript-specific algebra checks and reproducibility}
The expanded sources, Hbar covectors, pullbacks, ground Sigma/Pi reduction,
normalization multipliers and Euler identities are checked by
\path{devtools/check_ccpq_theory.py}. It compares them with the independent
complex-step oracle for closed/nonclosed P, ground/first two excited states,
and both denominator models: 12 configurations. The largest discrepancy
in the saved run is $\ee{8.882}{-16}$. In the nonclosed examples the script
also asserts $J\ne B$ and $\lambda\ne\lambda^g$.

The audit deliberately keeps spin partners together in its small nonclosed
P so this machine-precision formula check is not dominated by a near-zero
Q denominator. Existing stress tests separately exercise difficult
corrections; they do not turn small-denominator models into reliable physical
predictions. The broader gradient regression run contains 144 passing tests;
these 12 manuscript algebra configurations are an additional check, not a
relabeling of that pytest count.

Eight additional fixed-Fock reduced-response checks are described in
Sec.~\ref{sec:leanvalidation}; their scope and count are separate.

Run from the repository root with the CCpy environment active:
\begin{verbatim}
OPENBLAS_NUM_THREADS=1 python devtools/check_ccpq_theory.py
OPENBLAS_NUM_THREADS=1 python devtools/check_ccpq_lean_response.py
OPENBLAS_NUM_THREADS=1 python -m pytest tests/gradients \
    tests/scf/test_ccpq_gradient.py -q
\end{verbatim}
For molecular reproductions, build the optional FCIpy connector using
\path{devtools/build_fcipy_oracle.py} and set its printed
\texttt{CCPY\_FCIPY\_EXTENSION}. The optimized comparison examples are
\path{examples/ccpq_optimized_comparison.py} and
\path{examples/ccpq_excited_optimized_comparison.py}.
General substantial-T/Q demonstrations use
\path{examples/ccpq_complete_spaces.py}; exact command lines and fixed
selection lists are in the result reports named below.

\begingroup\small\setlength{\parskip}{0pt}
\subsection*{Saved evidence index}
\path{docs/theory/ccpq_validation_manifest.json} records SHA-256 hashes.
Repository-relative paths and hashes record provenance, not correctness.

Reports in \path{examples/results/}:
\par\noindent\path{h2o_631g_optimized_df.json}\quad\path{h2o_631g_excited_optimized.json}
\par\noindent\path{h4_631g_excited_optimized.json}\quad\path{h4_631g_excited_step_convergence.json}
\par\noindent\path{h6_healthy_v6.json}\quad\path{h6_healthy_v7.json}
\par\noindent\path{h2o_631g_frozen_core.json}\quad\path{h2o_631g_reduced_all_components.json}
\par Checks in \path{docs/theory/}:
\par\noindent\path{ccpq_formula_checks.json}
\par\noindent\path{ccpq_lean_response_checks.json}
\par\noindent\path{ccpq_regression_run.txt}
\par Tables: \path{devtools/build_ccpq_theory_evidence.py}.
\par Master: \path{docs/theory/ccpq_analytical_gradients.tex}.
\endgroup


\end{document}
